\documentclass[a4paper,twoside]{article}

\usepackage{morewrites}

\usepackage[svgnames,dvipsnames,x11names]{xcolor}
\usepackage{graphicx}
\usepackage[english]{babel}
\usepackage[colorlinks=true, linkcolor=NavyBlue, urlcolor=NavyBlue, citecolor=NavyBlue]{hyperref}
\usepackage[a4paper]{geometry}
\usepackage{ts-fonts}
\usepackage{amsmath}
\usepackage{float}
\usepackage{seqsplit}
\usepackage{ts-callouts}
\usepackage{ts-code}

\usepackage{lastpage}
\usepackage{refcount}
\makeatletter
\def\ps@plain{
  \let\@oddhead\@empty
  \let\@evenhead\@empty
  \def\@oddfoot{
    \hfil
    \ifnum\getpagerefnumber{LastPage}>1\relax
      \thepage
    \fi
    \hfil}
  \let\@evenfoot\@oddfoot
}
\makeatother

\title{Getting Started}
\author{}\date{}

\begin{document}

\maketitle

\thispagestyle{plain}
In our journey to typeset beautiful documents with TeXSmith, we'll start with the basics:

\begin{enumerate}
\item Turn Markdown or HTML into LaTeX/PDF.
\item Drop TeXSmith into an existing MkDocs site.
\item Drive it from Python.

\end{enumerate}
Hop to the section you need or read straight through for the big picture.

\section{Installation}\label{installation}

To install TeXSmith, use your preferred Python package manager:

\textbf{pip}\par

\begin{code}{bash}{}{}
\begin{Verbatim}[commandchars=\\\{\},breaklines, breakanywhere, commandchars=\\\{\}]
pip\PY{+w}{ }install\PY{+w}{ }texsmith
\end{Verbatim}
\end{code}
\textbf{pipx}\par

\begin{code}{bash}{}{}
\begin{Verbatim}[commandchars=\\\{\},breaklines, breakanywhere, commandchars=\\\{\}]
pipx\PY{+w}{ }install\PY{+w}{ }texsmith
\end{Verbatim}
\end{code}
\textbf{uv}\par

\begin{code}{bash}{}{}
\begin{Verbatim}[commandchars=\\\{\},breaklines, breakanywhere, commandchars=\\\{\}]
uv\PY{+w}{ }tool\PY{+w}{ }install\PY{+w}{ }texsmith
\end{Verbatim}
\end{code}
For basic use, you don't need anything else. TeXSmith bundles Tectonic for \LaTeX{} builds and will auto-install the required tools on demand.

To build with the \textbf{Typst} backend (\texttt{-\allowbreak{}-\allowbreak{}format typst}), install the embedded compiler as an extra so no system binary is required:

\textbf{pip}\par

\begin{code}{bash}{}{}
\begin{Verbatim}[commandchars=\\\{\},breaklines, breakanywhere, commandchars=\\\{\}]
pip\PY{+w}{ }install\PY{+w}{ }\PY{l+s+s2}{\PYZdq{}texsmith[typst]\PYZdq{}}
\end{Verbatim}
\end{code}
\textbf{uv}\par

\begin{code}{bash}{}{}
\begin{Verbatim}[commandchars=\\\{\},breaklines, breakanywhere, commandchars=\\\{\}]
uv\PY{+w}{ }tool\PY{+w}{ }install\PY{+w}{ }\PY{l+s+s2}{\PYZdq{}texsmith[typst]\PYZdq{}}
\end{Verbatim}
\end{code}
A system \texttt{typst} binary on your \texttt{PATH} is detected automatically as a fallback. See \hyperref[snippet:<HASH>]{Output backends} for details and the math caveat.

\section{Convert a Markdown file to LaTeX}\label{convert-a-markdown-file-to-latex}

By default TeXSmith writes \LaTeX{} to stdout. Pipe it or direct it into a folder. HTML works too:

\textbf{Here document}\par

\begin{code}{text}{}{}
\begin{Verbatim}[commandchars=\\\{\},breaklines, breakanywhere, commandchars=\\\{\}]
cat \PYZlt{}\PYZlt{} EOF | texsmith
\PYZsh{} Title

Some **bold** text.

\PYZhy{} Foo
\PYZhy{} Bar
EOF
\PYZbs{}section\PYZob{}Title\PYZcb{}\PYZbs{}label\PYZob{}title\PYZcb{}

Some \PYZbs{}textbf\PYZob{}bold\PYZcb{} text.

\PYZbs{}begin\PYZob{}itemize\PYZcb{}
\PYZbs{}item\PYZob{}\PYZcb{} Foo
\PYZbs{}item\PYZob{}\PYZcb{} Bar

\PYZbs{}end\PYZob{}itemize\PYZcb{}
\end{Verbatim}
\end{code}
\textbf{From file}\par

\begin{code}{text}{}{}
\begin{Verbatim}[commandchars=\\\{\},breaklines, breakanywhere, commandchars=\\\{\}]
\PYZdl{} echo \PYZdq{}\PYZsh{} Title\PYZbs{}nSome **bold** text.\PYZdq{} \PYZgt{} sample.md
\PYZdl{} texsmith sample.md \PYZhy{}\PYZhy{}output build/
\PYZbs{}chapter\PYZob{}Title\PYZcb{}\PYZbs{}label\PYZob{}title\PYZcb{}

Some \PYZbs{}textbf\PYZob{}bold\PYZcb{} text.
\end{Verbatim}
\end{code}
\textbf{HTML}\par

\begin{code}{text}{}{}
\begin{Verbatim}[commandchars=\\\{\},breaklines, breakanywhere, commandchars=\\\{\}]
\PYZdl{} echo \PYZdq{}\PYZlt{}h1\PYZgt{}Title\PYZlt{}/h1\PYZgt{}\PYZlt{}p\PYZgt{}Some \PYZlt{}strong\PYZgt{}bold\PYZlt{}/strong\PYZgt{} text.\PYZlt{}/p\PYZgt{}\PYZdq{} \PYZgt{} sample.html
\PYZdl{} texsmith sample.html
\PYZbs{}chapter\PYZob{}Title\PYZcb{}
Some \PYZbs{}textbf\PYZob{}bold\PYZcb{} text.
\end{Verbatim}
\end{code}
\section{Generate a PDF}\label{generate-a-pdf}

Want the full PDF? Start with our playful \href{https://en.wikipedia.org/wiki/Booby}{booby} example or create your own \texttt{booby.md}:

\begin{code}{markdown}{}{}
\begin{Verbatim}[commandchars=\\\{\},breaklines, breakanywhere, commandchars=\\\{\}]

\end{Verbatim}
\end{code}
Notice the front matter up top: it carries the title, author, date, and template to use.

Then let TeXSmith crunch it:

\begin{code}{bash}{}{}
\begin{Verbatim}[commandchars=\\\{\},breaklines, breakanywhere, commandchars=\\\{\}]
texsmith\PY{+w}{ }booby.md\PY{+w}{ }\PYZhy{}\PYZhy{}output\PY{+w}{ }build/\PY{+w}{ }\PYZhy{}apaper\PY{o}{=}a5\PY{+w}{ }\PYZhy{}\PYZhy{}build
\end{Verbatim}
\end{code}
With Tectonic as the default engine, fonts, packages, and dependencies resolve themselves on demand (including Tectonic if it is missing). Nothing else to install.

Enjoy a fresh PDF at \texttt{build/booby.pdf}:

\begin{figure}[H]
    \centering

    \ifdefined\adjustbox

        \adjustbox{max width=\textwidth}{
            \includegraphics[width=0.7\linewidth]{snippet-<HASH>.pdf}
        }

    \else

        \includegraphics[width=0.7\linewidth]{snippet-<HASH>.pdf}

    \fi

\end{figure}
Peek inside \texttt{build/} to find \texttt{booby.tex}. Swap \texttt{-\allowbreak{}-\allowbreak{}template} when you want a different \LaTeX{} project layout or polish level:

\begin{code}{bash}{}{}
\begin{Verbatim}[commandchars=\\\{\},breaklines, breakanywhere, commandchars=\\\{\}]
texsmith\PY{+w}{ }booby.md\PY{+w}{ }\PYZhy{}\PYZhy{}template\PY{+w}{ }article\PY{+w}{ }\PYZhy{}\PYZhy{}output\PYZhy{}dir\PY{+w}{ }build
\end{Verbatim}
\end{code}
The default toolchain is \texttt{tectonic}, which auto-installs itself and required packages. If you prefer using your system \LaTeX{} installation, specify \texttt{-\allowbreak{}-\allowbreak{}engine lualatex} or \texttt{-\allowbreak{}-\allowbreak{}engine xelatex} instead. Both commands yield \texttt{doc.pdf} in the current directory. Open it to see the rendered output.

If you want to customize the layout, choose a template with \texttt{-\allowbreak{}-\allowbreak{}template article}, \texttt{-\allowbreak{}-\allowbreak{}template book} or \texttt{-\allowbreak{}-\allowbreak{}template your-\allowbreak{}own-\allowbreak{}template}.

You may want to pass additional \LaTeX{} options such as \texttt{-\allowbreak{}apaper=a4} or \texttt{-\allowbreak{}amargin=1in} to tweak page geometry:

\section{Optional prerequisites}\label{optional-prerequisites}

\begin{description}
\item[{ LaTeX distribution }] Install TeX Live, MiKTeX, or MacTeX if you want TeXSmith to hand off builds to \texttt{latexmk} (\texttt{-\allowbreak{}-\allowbreak{}engine lualatex} / \texttt{-\allowbreak{}-\allowbreak{}engine xelatex}). The default route uses Tectonic, which auto-installs itself and required packages.
\item[{ Typst compiler }] Needed only for \texttt{-\allowbreak{}-\allowbreak{}format typst -\allowbreak{}-\allowbreak{}build}. Install the embedded compiler with \texttt{pip install "texsmith[typst]"}, or put a \texttt{typst} binary on your \texttt{PATH}. Emitting the \texttt{.typ} source (without \texttt{-\allowbreak{}-\allowbreak{}build}) needs no compiler. See \hyperref[snippet:<HASH>]{Output backends}.
\item[{ Diagram tooling }] Mermaid-to-PDF (\texttt{minlag/mermaid-\allowbreak{}cli}) conversion falls back to Docker. Install Docker Desktop (with WSL integration on Windows) or register your own converter if Mermaid diagrams are common in your docs.

\end{description}
Draw.io and Mermaid diagrams try a Playwright exporter first (cached under \texttt{\textasciitilde{}/.cache/texsmith/playwright}), then the local CLI, then Docker (\texttt{rlespinasse/drawio-\allowbreak{}desktop-\allowbreak{}headless} / \texttt{minlag/mermaid-\allowbreak{}cli}). Use \texttt{-\allowbreak{}-\allowbreak{}diagrams-\allowbreak{}backend=playwright|local|docker} to pin a specific backend.

\begin{description}
\item[{ Fonts }] TeXSmith ships with Noto fallback for wide Unicode coverage. Add your own fonts if you want a specific script or branded look.
\item[{ Legacy LaTeX accents }] By default TeXSmith emits Unicode glyphs. If you need legacy LaTeX accent macros, pass \texttt{-\allowbreak{}-\allowbreak{}legacy-\allowbreak{}latex-\allowbreak{}accents} on the CLI or set \texttt{ConversionRequest(legacy\_latex\_accents=True)} in the API.

\end{description}
\section{Use the Python API}\label{use-the-python-api}

TeXSmith also ships as a Python library. Create \texttt{demo.py}:

\begin{code}{python}{}{}
\begin{Verbatim}[commandchars=\\\{\},breaklines, breakanywhere, commandchars=\\\{\}]
\PY{k+kn}{from}\PY{+w}{ }\PY{n+nn}{pathlib}\PY{+w}{ }\PY{k+kn}{import} \PY{n}{Path}
\PY{k+kn}{from}\PY{+w}{ }\PY{n+nn}{texsmith}\PY{+w}{ }\PY{k+kn}{import} \PY{n}{Document}\PY{p}{,} \PY{n}{convert\PYZus{}documents}

\PY{n}{bundle} \PY{o}{=} \PY{n}{convert\PYZus{}documents}\PY{p}{(}
    \PY{p}{[}\PY{n}{Document}\PY{o}{.}\PY{n}{from\PYZus{}markdown}\PY{p}{(}\PY{n}{Path}\PY{p}{(}\PY{l+s+s2}{\PYZdq{}}\PY{l+s+s2}{intro.md}\PY{l+s+s2}{\PYZdq{}}\PY{p}{)}\PY{p}{)}\PY{p}{]}\PY{p}{,}
    \PY{n}{output\PYZus{}dir}\PY{o}{=}\PY{n}{Path}\PY{p}{(}\PY{l+s+s2}{\PYZdq{}}\PY{l+s+s2}{build}\PY{l+s+s2}{\PYZdq{}}\PY{p}{)}\PY{p}{,}
\PY{p}{)}

\PY{n+nb}{print}\PY{p}{(}\PY{l+s+s2}{\PYZdq{}}\PY{l+s+s2}{Fragments:}\PY{l+s+s2}{\PYZdq{}}\PY{p}{,} \PY{p}{[}\PY{n}{fragment}\PY{o}{.}\PY{n}{stem} \PY{k}{for} \PY{n}{fragment} \PY{o+ow}{in} \PY{n}{bundle}\PY{o}{.}\PY{n}{fragments}\PY{p}{]}\PY{p}{)}
\PY{n+nb}{print}\PY{p}{(}\PY{l+s+s2}{\PYZdq{}}\PY{l+s+s2}{Preview:}\PY{l+s+s2}{\PYZdq{}}\PY{p}{,} \PY{n}{bundle}\PY{o}{.}\PY{n}{combined\PYZus{}output}\PY{p}{(}\PY{p}{)}\PY{p}{[}\PY{p}{:}\PY{l+m+mi}{120}\PY{p}{]}\PY{p}{)}
\end{Verbatim}
\end{code}
Run the snippet with \texttt{python demo.py}. The API mirrors the CLI; reach for \texttt{ConversionService} or \texttt{TemplateSession} when you need fine-grained control over slot assignments, diagnostics, or template metadata.

\section{Convert a MkDocs site}\label{convert-a-mkdocs-site}

Point TeXSmith at a MkDocs site after \texttt{mkdocs build} renders clean HTML:

\begin{code}{bash}{}{}
\begin{Verbatim}[commandchars=\\\{\},breaklines, breakanywhere, commandchars=\\\{\}]
\PY{c+c1}{\PYZsh{} Build your MkDocs site into a disposable directory}
mkdocs\PY{+w}{ }build

\PY{c+c1}{\PYZsh{} Convert one page into LaTeX/PDF\PYZhy{}ready assets}
texsmith\PY{+w}{ }build/site/guides/overview/index.html\PY{+w}{ }\PY{l+s+se}{\PYZbs{}}
\PY{+w}{  }\PYZhy{}\PYZhy{}template\PY{+w}{ }article\PY{+w}{ }\PY{l+s+se}{\PYZbs{}}
\PY{+w}{  }\PYZhy{}\PYZhy{}output\PYZhy{}dir\PY{+w}{ }build/press\PY{+w}{ }\PY{l+s+se}{\PYZbs{}}
\PY{+w}{  }docs/references.bib
\end{Verbatim}
\end{code}
\begin{callout}[callout tip]{Tip}
The default selector (\texttt{article.md-\allowbreak{}content\_\_inner}) already matches MkDocs Material content; skip \texttt{-\allowbreak{}-\allowbreak{}selector} unless you heavily customise templates.

When your site spans multiple documents, repeat the command per page and stitch them together with template slots (for example, \texttt{-\allowbreak{}-\allowbreak{}slot mainmatter:build/site/manual/index.html}).

For live previews, point TeXSmith at the temporary site directory that \texttt{mkdocs serve} prints on startup.

Once the \LaTeX{} bundle looks good, add \texttt{-\allowbreak{}-\allowbreak{}build} to invoke your engine of choice or wire it into CI so MkDocs HTML \(\rightarrow\) TeXSmith PDF runs on every build.
\end{callout}

\end{document}
